\label{sec:boundary_conditions}

Le domaine de calcul est un rectangle fermé
$\Omega = [0,L_x] \times [0,L_y]$ représenté par les bords de notre simulation, délimité par quatre parois fixes.
Le vecteur $\mathbf{n}$ désigne la normale unitaire sortante à la frontière
$\partial\Omega$.
Les conditions aux limites précisent le comportement de la vitesse,
de la pression et du colorant sur ces parois.

\paragraph*{Vitesse : parois imperméables}

Le fluide ne traverse pas les parois : la composante normale de la vitesse
est nulle sur la frontière.
\begin{equation}
    \label{eq:wall_impermeability}
    \boxed{\mathbf{u} \cdot \mathbf{n} = 0
        \quad \text{sur } \partial\Omega.}
\end{equation}
Avec $\mathbf{u} = (u,v)$, cette condition s'écrit :
\begin{equation}
    \begin{aligned}
        u(0,y,t) = u(L_x,y,t) &= 0, \\
        v(x,0,t) = v(x,L_y,t) &= 0.
    \end{aligned}
\end{equation}

La composante tangentielle n'est pas imposée à zéro : le fluide peut
glisser le long des murs.
Cette condition convient au modèle non visqueux actuellement implémenté.
Les valeurs tangentielles utilisées au bord sont prolongées par symétrie,
sans imposer d'adhérence à la paroi.

Pour une extension visqueuse, le comportement tangentiel devra être
précisé : une paroi fixe avec adhérence impose $\mathbf{u}=\mathbf{0}$,
tandis qu'une paroi à glissement libre impose une contrainte visqueuse
tangentielle nulle.
Sur une paroi plane, avec une viscosité constante non nulle, ce dernier
choix correspond à $\partial u_{\mathrm{t}}/\partial n = 0$,
où $u_{\mathrm{t}}$ est la composante tangentielle de la vitesse.

\paragraph*{Colorant : absence de flux aux parois}

En l'absence de diffusion, le flux de colorant est $c\mathbf{u}$.
L'imperméabilité des murs entraîne donc :
\begin{equation}
    \label{eq:dye_boundary}
    \boxed{(c\mathbf{u})\cdot\mathbf{n} = 0
        \quad \text{sur } \partial\Omega.}
\end{equation}
Aucune concentration nulle n'est imposée au bord : le colorant peut
être présent contre une paroi, mais ne peut pas la traverser.
Dans ce domaine fermé, sa quantité totale ne varie que sous l'effet
du terme source $s$ :
\begin{equation}
    \frac{\mathrm{d}}{\mathrm{d}t}\int_{\Omega} c\,\mathrm{d}\Omega
    = \int_{\Omega} s\,\mathrm{d}\Omega.
\end{equation}

\paragraph*{Conditions initiales}

La configuration initiale de KIRIN correspond à un fluide au repos
et à un domaine sans colorant :
\begin{equation}
    \mathbf{u}(\mathbf{x},0) = \mathbf{0},
    \qquad c(\mathbf{x},0) = 0.
\end{equation}
La pression est initialisée à zéro comme point de départ du solveur
itératif ; elle est ensuite déterminée par la contrainte
d'incompressibilité.
Les interactions de l'utilisateur ajoutent du colorant et modifient
la vitesse à l'intérieur du domaine, sans autoriser de flux à travers
les murs.
